\chapter{Discussion}
\label{ch:discussion}

Les chapitres précédents ont décrit la pipeline telle qu'elle a été construite,
calibrée et exécutée. Ce chapitre prend du recul. Plutôt que de présenter une
liste de fonctionnalités, je reviens sur ce que la mise en {\oe}uvre a réellement
révélé : les limites que chaque module a fini par exposer, les choix techniques
qui se sont imposés une fois confrontés au matériel, et ce que la confrontation
entre une approche modulaire et le paradigme de l'imitation laisse entrevoir.
Plusieurs de ces points ne sont pas des décisions prises \emph{a priori} : ce
sont des conclusions tirées d'essais qui ont d'abord échoué. Je m'efforce de les
restituer dans cet ordre --- hypothèse, essai, problème observé, correctif ---
car c'est là que se trouve l'essentiel de l'apprentissage du projet.

\section{Synthèse des limites identifiées}
\label{sec:disc_limites}

Maintenant que la pipeline complète est posée, il devient possible d'en cartographier
les points faibles de façon ordonnée. Chaque limite ci-dessous a été rencontrée
en pratique, comprise, puis --- quand cela était possible --- partiellement
corrigée. Identifier honnêtement ce qui ne marche pas est, à mes yeux, aussi
important que documenter ce qui marche.

\subsection{Limites de la perception V1 (HSV)}

Le détecteur HSV de première version partait d'une hypothèse simple : un objet
de couleur saturée se distingue de son fond par sa teinte. Cette hypothèse tient
dans un décor contrôlé, mais s'effondre dès que l'environnement contient des
surfaces de teinte voisine. Les limites observées sont irréductibles au seul
seuillage colorimétrique :
\begin{itemize}
  \item Confusion couleur entre l'objet cible et l'environnement : le cube orange
        est confondu avec le filament orange du robot (même teinte $H$), le mug
        blanc avec le sol blanc. Le blanc et le noir ne sont d'ailleurs pas des
        couleurs au sens HSV --- une luminosité $V$ quasi nulle rend $H$ et $S$
        indéterminés, une saturation $S$ quasi nulle rend $H$ indéterminé
        (\cite{forsyth_computer_2012}).
  \item Aucune notion de forme ni de contexte : seul le pixel compte.
  \item Sensibilité aux variations d'éclairage, qui obligent à recalibrer les
        plages HSV.
\end{itemize}

Ces limites ont été partiellement palliées par des zones d'exclusion géométriques
(un cylindre autour de la base, des bornes de workspace dans \texttt{scene.json}).
Mais elles empêchaient une validation fiable de l'aval de la pipeline : on ne
pouvait pas distinguer une erreur de planification d'une erreur de détection.
\textbf{Correctif} : le passage au V2 (\texttt{HFDetector}), initialement prévu en
fin de projet, a été \emph{anticipé} précisément pour cette raison.

\subsection{Limites de la perception V2 (OWL-ViTv2)}

Le V2 résout le problème de fond du V1 --- il raisonne sur une description
textuelle enrichie (couleur, forme, matière, trait distinctif) plutôt que sur la
seule teinte --- mais il introduit ses propres faiblesses, de nature sémantique
cette fois :
\begin{itemize}
  \item Faux positifs sémantiques : une gourde en aluminium bleu marine détectée
        comme \texttt{purple\_cylinder}, un essuie-tout comme \texttt{white\_mug}.
        Le modèle se trompe non plus sur le pixel mais sur le concept.
  \item Vitesse limitée ($\sim$3--5~FPS sur le MacBook M4 en MPS), trop lente pour
        un asservissement haute fréquence.
  \item Pas de masque de segmentation : seulement des \emph{bounding boxes}, ce qui
        prive le grasp planner d'un contour précis.
\end{itemize}

Une limite supplémentaire est apparue tardivement, sur le cylindre violet utilisé
comme objet de test principal : sa \emph{détection HSV de secours} est très
faible, le score d'aire tombant à son plancher ($\sim$0{,}05). Ce score faible
dégrade à la fois la position triangulée et l'orientation estimée du grand axe. Recalibrer le
descripteur textuel ou les plages HSV de secours améliorerait les deux
simultanément --- c'est une limite documentée, non un défaut de conception.

\textbf{Pistes de correctif} : descriptions plus précises, filtrage par zone de
travail (déjà en place), seuil de confiance plus strict, mode hybride HSV (rapide)
+ HF (validation périodique), et l'ajout d'un masque par SAM~2 en post-traitement
pour récupérer un contour.

\subsection{Limites du grasp planner (top-down)}

La planification de saisie réellement implémentée est une heuristique top-down
(\cite{bohg_data-driven_2014}) : la pince pointe vers le bas et le \emph{yaw}
(\texttt{wrist\_roll}) reste libre. Cette stratégie est robuste, mais elle borne
de fait l'ensemble des objets saisissables :
\begin{itemize}
  \item Un objet trop haut ($>12$~cm) est rejeté : la pince ne peut pas
        l'approcher par le dessus sans collision de poignet.
  \item Un objet fin demande un alignement \texttt{wrist\_roll} très précis, à la
        limite de ce que permet le plancher de bruit de répétabilité.
  \item La pince TPU est compliante et le cylindre de test est rond : si le centre
        de prise est légèrement décentré, l'objet peut glisser entre les mâchoires
        au moment du serrage --- un risque de glissement qu'une campagne formelle
        permettrait de quantifier.
\end{itemize}

La principale limite n'est cependant pas le top-down lui-même mais le fait qu'il
soit la \emph{seule} approche fiable disponible sur ce bras. La section
\ref{sec:disc_choix} y revient en détail, car c'est un résultat à part entière.

\subsection{Limites de la précision géométrique}

Les résidus \emph{hand-eye} mesurés --- $\sim$6{,}6~mm en moyenne pour
\texttt{cam\_0}, $\sim$11{,}8~mm pour \texttt{cam\_1} --- restent dans l'ordre de
grandeur du \textbf{bruit de répétabilité du SO-101}
($\sim$5--10~mm en spécification constructeur) ; la caméra \texttt{cam\_2}
eye-in-hand fait nettement mieux (2{,}48/4{,}77~mm) grâce à un montage plus rigide. Cette
imprécision se propage en aval :
\begin{itemize}
  \item L'objet est localisé en 3D à $\pm$5--10~mm près.
  \item Le grasp planner positionne la pince à la même précision.
  \item L'ouverture adaptative de la pince laisse une marge de 10~mm de chaque
        côté de l'objet, ce qui absorbe cette incertitude pour les objets
        non-fins mais reste tendu pour un stylo.
\end{itemize}

\textbf{Correctifs mis en {\oe}uvre} : d'abord, ancrer la profondeur de prise sur
le plan de la table ($z_\text{grasp} = z_\text{table} + h_\text{objet}/2$) plutôt
que sur le $Z$ bruité de la stéréo, dont l'erreur croît en $Z^2$ ; ensuite, la
replanification en boucle fermée, où \texttt{cam\_2} confirme la position --- et
désormais l'orientation --- de l'objet à la pose d'approche et corrige la
trajectoire par intersection rayon-plan avant la descente finale.

\section{Choix académiques discutables et résultat sur la sous-actionation}
\label{sec:disc_choix}

Toute limite de perception peut, en principe, se corriger par un meilleur modèle.
La sous-actionation du bras, elle, est structurelle : c'est une contrainte
mécanique que ni le code ni la calibration ne peuvent contourner. Cette section
explicite trois choix qui méritent d'être justifiés, dont le dernier --- le
renoncement à la prise latérale --- découle directement de la cinématique du
bras.

\subsection{Pourquoi OWL-ViTv2 plutôt que Grounding-DINO ?}

\todoF{Voix de Maxence à confirmer. Synthèse : OWL-ViTv2
(\cite{minderer_owlv2_2023}) est mieux documenté et plus standardisé dans la
stack Hugging Face --- la même librairie \texttt{transformers} que les politiques
LeRobot ---, ce qui réduit la dette d'intégration. Grounding-DINO
(\cite{liu_grounding_dino_2024}) serait sans doute plus précis sur les objets
fins, mais au prix d'une dépendance supplémentaire et d'une intégration moins
directe. Le choix penche vers la cohérence de l'écosystème.}

\subsection{Pourquoi pas un détecteur fine-tuné (YOLO) ?}

\todoF{Voix de Maxence à confirmer. Synthèse : un détecteur fine-tuné aurait
demandé un jeu d'images annotées par objet, soit plusieurs jours de collecte et
d'annotation. La détection open-vocabulary élimine ce coût au prix d'une part de
précision --- un compromis raisonnable pour un travail de bachelor où l'effort
porte sur l'orchestration de la pipeline, pas sur l'optimisation d'un détecteur.}

\subsection{Pourquoi 5 DDL, et ce que cela implique}

Le SO-101 est un robot \textbf{sous-actionné} : il possède cinq articulations
rotoïdes actionnées pour un effecteur qui vit dans $SE(3)$, espace à six degrés de
liberté. Conséquence directe : tous les couples (position, orientation) ne sont
pas atteignables. Ce n'est pas un défaut à corriger mais une donnée du problème,
qu'il faut comprendre pour décider quelles saisies tenter.

Pour la saisie top-down, on contraint l'orientation à un seul degré de liberté
libre (\texttt{wrist\_roll}) : la pose désirée fixe la position (3 contraintes) et
deux composantes de l'orientation (axe d'approche vertical), tandis que le
\emph{yaw} reste libre. On a donc cinq contraintes pour cinq articulations, dont
une laissée libre et compensée par l'IK --- \textbf{le système est bien posé}, et
le roundtrip FK$\rightarrow$IK le confirme (erreur $\sim$0{,}8~mm sur des
configurations aléatoires) ; une pose top-down typique se résout à $\sim$9{,}7~mm,
résidu de la sous-actuation que le solveur signale comme non exactement convergé.

\paragraph{Ce que le bras peut et ne peut pas faire.}
Les trois articulations de tangage (\texttt{shoulder\_lift}, \texttt{elbow\_flex},
\texttt{wrist\_flex}) plient toutes dans un \emph{même} plan vertical. Le bras
atteint donc un large éventail d'orientations d'approche \emph{à l'intérieur} de
ce plan --- verticale, inclinée, voire horizontale frontale --- mais il ne peut
pas orienter l'approche \emph{hors} de ce plan. C'est cette restriction qui
détermine ce qui est préhensible.

\paragraph{Pourquoi la prise latérale n'a pas été retenue.}
La question s'est posée concrètement : pourrait-on saisir un objet par le flanc,
pince horizontale se refermant autour de sa paroi verticale, comme une main qui
attrape une canette ? Une telle fermeture exige précisément d'orienter l'approche
\emph{hors} du plan de tangage décrit ci-dessus --- ce qui, sur ce bras à cinq
articulations, n'est pas atteignable de façon fiable. La prise latérale n'a donc
\textbf{pas été implémentée} : le travail s'en tient à la saisie top-down. Sans
être catégoriquement impossible, cette famille de prises n'est pas robuste sur
5~DDL ; un bras à six degrés de liberté lèverait la contrainte en rendant
l'orientation d'approche libre. Ce choix borne mécaniquement l'ensemble des objets
saisissables : un objet trop haut, ou en position debout hors de portée par le
dessus, sort de l'espace préhensible.

\subsection{La convention de fermeture à 90\textdegree{} : un débogage révélateur}

Un dernier choix mérite d'être raconté comme tel, parce qu'il illustre le décalage
entre le modèle et le réel. L'orientation de prise calculée par le grasp planner
visait à aligner les mâchoires en travers du petit axe de l'empreinte de l'objet.
En théorie, l'angle envoyé au \texttt{wrist\_roll} devait suffire. En pratique, les
saisies échouaient systématiquement : la pince se présentait toujours à
contresens, comme tournée d'un quart de tour.

L'hypothèse initiale --- une erreur de signe dans la rotation $R(\theta)$ --- ne
résistait pas à la vérification : le déterminant valait bien $+1$ et la FK
recomposait la pose visée. Le problème ne venait donc pas du calcul mais de la
\emph{convention} reliant l'angle du modèle au montage physique des mâchoires. Pour
le confirmer, j'ai ajouté un log de diagnostic affichant l'angle vu par chaque
caméra (\texttt{yaw\_cam0\_deg}, \texttt{yaw\_cam1\_deg}) et exposé une option de
test (\texttt{grasp-yaw-offset~90}). L'essai a été sans appel : un décalage de
$90$\textdegree{} faisait converger l'alignement. Les mâchoires se ferment donc à
$90$\textdegree{} de la convention nominale du code --- un fait qui n'apparaît dans
aucun document, et qu'aucune dérivation mathématique ne pouvait révéler. La valeur
par défaut de déploiement a été fixée à $90$\textdegree{}. Cet épisode résume bien
la nature du travail : une partie des difficultés ne se résout ni par les
mathématiques ni par la calibration, mais par l'instrumentation patiente de
l'écart entre le code et la mécanique.

\section{Comparaison des deux paradigmes (modulaire vs imitation)}
\label{sec:disc_paradigmes}

La pipeline construite ici relève d'une approche \emph{modulaire} : perception,
planification et contrôle sont des modules distincts, chacun analysable et
remplaçable. Le paradigme concurrent --- l'imitation learning, dont SmolVLA
(\cite{shukor_smolvla_2025}) est un représentant accessible sur ce même matériel
--- apprend de bout en bout une politique sensorimotrice à partir de
démonstrations. Comparer les deux est l'un des objectifs de mise en perspective de
ce travail. Je le fais ici sur le plan \emph{qualitatif}, en distinguant clairement
ce qui est argumenté de ce qui resterait à mesurer.

\paragraph{Avantages de l'architecture modulaire.}
\begin{itemize}
  \item Explicabilité : en cas d'échec, on peut isoler l'étape fautive (détection,
        triangulation, IK, contrôle moteur).
  \item Modularité : on remplace un module sans toucher au reste --- le passage du
        détecteur V1 au V2 en est l'illustration directe.
  \item Garanties géométriques : on \emph{sait} où la pince est censée se trouver,
        car la position est calculée explicitement par l'IK et vérifiable par FK.
  \item Évitement d'obstacles \emph{explicite} : une trajectoire peut être vérifiée
        contre une représentation de la scène avant exécution.
\end{itemize}

\paragraph{Avantages de l'imitation learning.}
\begin{itemize}
  \item Pas de \emph{feature engineering} : le réseau apprend ses propres
        représentations à partir des images et des états.
  \item Robustesse potentielle à des situations non anticipées, dans la mesure où
        elles sont représentées dans le jeu de démonstrations.
  \item Code de déploiement minimal : un appel à \texttt{policy.select\_action}
        remplace toute l'orchestration décrite aux chapitres précédents.
\end{itemize}

\paragraph{Inconvénients respectifs.}
\begin{itemize}
  \item Modulaire : forte charge d'ingénierie ; chaque module est un point de
        défaillance, et la décision à 90\textdegree{} ci-dessus montre que les
        conventions implicites coûtent cher.
  \item Imitation : boîte noire (un échec est difficile à expliquer) ; dépendance
        à un jeu de démonstrations de qualité ; aucune garantie formelle
        d'évitement de collision.
\end{itemize}

\todoF{La comparaison \emph{empirique} entre les deux paradigmes n'a PAS été
réalisée. Le banc V1/V2 et une politique SmolVLA n'ont pas été évalués côte à côte
sur un protocole commun (taux de réussite, nombre de replans, collisions évitées,
précision de prise). Cette comparaison est donc présentée ici comme une
\emph{perspective}, pas comme une contribution réalisée. Les tableaux chiffrés
correspondants (chapitre Résultats) restent des placeholders tant que la campagne
formelle --- script \texttt{experiment\_campaign.py} --- n'a pas été menée ni
journalisée.}

\section{Pistes d'amélioration}
\label{sec:disc_pistes}

Cette mise en perspective ouvre naturellement sur les directions qui prolongeraient
le travail. Plusieurs découlent directement des limites ci-dessus :
\begin{itemize}
  \item \textbf{Segmentation} : ajouter SAM~2 en post-traitement pour fournir au
        grasp planner un contour, et non plus une simple \emph{bounding box} ---
        ce qui affinerait l'estimation d'orientation aujourd'hui dépendante d'une
        détection au score parfois faible.
  \item \textbf{Planification de trajectoire} : remplacer l'enchaînement de
        sous-trajectoires quintiques par un planificateur d'échantillonnage de type
        RRT (\cite{lavalle_planning_2006}) pour un évitement d'obstacles réellement
        général.
  \item \textbf{Active vision} : sélectionner dynamiquement le point de vue le moins
        occulté (objectif 2 du cahier des charges, aujourd'hui seulement effleuré
        par la coopération stéréo / eye-in-hand).
  \item \textbf{Asservissement visuel} : passer d'une correction discrète à la pose
        d'approche à un asservissement continu au sens de
        \cite{chaumette_visual_2006}.
  \item \textbf{Matériel} : un bras à six degrés de liberté lèverait la contrainte
        de sous-actionation et débloquerait les prises latérales aujourd'hui mal
        conditionnées.
  \item \textbf{Évaluation} : mener la campagne formelle et la comparaison
        empirique avec l'imitation learning, qui restent à ce stade des
        perspectives (cf. \ref{sec:disc_paradigmes}).
\end{itemize}

Ces pistes supposent toutes un socle commun : une pipeline complète, instrumentée
et reproductible. C'est ce socle, plus que tel ou tel résultat chiffré, que ce
travail aura cherché à établir. La conclusion qui suit en dresse le bilan d'ensemble
et le confronte aux six objectifs du cahier des charges.
